\par\medskip
\noindent
2.2. Nos proponemos precisar las relaciones entre $T_\ell(A_K)$ y
$T_\ell(A_o)=T_\ell(A_o^o)$ haciendo intervenir $T_\ell(A^o)$.
Recordemos primero lo siguiente.

\medskip
\noindent
\underline{Lema} 2.2.1. \underline{Sea $m>0$ un entero. Entonces son
equivalentes las condiciones siguientes:}

\smallskip
\noindent
(i) $m\,\mathrm{id}_{A^o}$ \underline{es plano}.

\noindent
(ii) $m\,\mathrm{id}_{A^o}$ \underline{es sobreyectivo}.

\noindent
(iii) $m\,\mathrm{id}_{A^o}$ \underline{es cuasifinito}.

\noindent
(iv) $m$ \underline{es primo con la característica de $k$}, o $A_o^o$
\underline{posee rango unipotente cero} (cf.
\sisi{1.5.5}{2.1.5}).

\smallskip
\noindent
\underline{Además, si estas condiciones se cumplen para $m$, se
cumplen para $m^\nu$ para todo $\nu>0$; los grupos}
\[
 {}_{m^\nu}A^o\underset{\mathrm{dfn}}{=}
 \ker(m^\nu\,\mathrm{id}_{A^o})
\]
\underline{son cuasifinitos, separados y planos sobre $S$, y, al variar
$\nu$, forman un sistema inverso de esquemas en grupos cuyos morfismos
de transición}
\[
 x\longmapsto m^{\nu'-\nu}x:
 {}_{m^{\nu'}}A^o\longrightarrow{}_{m^\nu}A^o
 \qquad(\nu'\geq\nu)
\]
\underline{son fielmente planos}.

En efecto, el criterio fibrado de planitud (EGA~IV~11.3.5) muestra que
(i), al igual que (ii) y (iii), es una condición sobre las fibras;
estas tres condiciones son equivalentes por SGA~3
$\mathrm{VI}_A$~5.6, puesto que las fibras en cuestión son lisas y
conexas. En la forma (ii), un teorema clásico de Weil [5, 15] muestra
que las condiciones se cumplen sobre la fibra genérica, que es un
esquema abeliano; basta, pues, enunciarlas sobre la fibra especial. Su
equivalencia con (iv) es entonces también bien conocida por la teoría
de estructura. En cualquiera de estas formas, la estabilidad de la
condición bajo $m\mapsto m^\nu$ es inmediata. Se sigue que los
${}_{m^\nu}A^o$ son esquemas en grupos cuasifinitos planos sobre $S$,
y son manifiestamente separados sobre $S$ porque $A^o$ lo es. Por
último, la planitud fiel de los morfismos de transición ---o,
equivalentemente, el hecho de que sean epimorfismos para la topología
fppf--- se deduce formalmente de que $m\,\mathrm{id}_{A^o}$ es
fielmente plano y, por tanto, un epimorfismo para dicha topología.

\medskip
\noindent
2.2.2. En consecuencia, tomaremos como $m$ un primo $\ell$ que
satisfaga las condiciones de 2.2.1. Denotamos entonces por
\begin{equation}\tag{2.2.2.1}
 T_\ell(A^o)=({}_{\ell^\nu}A^o)_{\nu>0}
\end{equation}
el sistema inverso descrito en 2.2.1. Mediante los dos cambios de base
$\eta\longrightarrow S$ y $s\longrightarrow S$, se recuperan los
sistemas inversos $T_\ell(A_K)$ y
$T_\ell(A_o^o)=T_\ell(A_o)$ considerados respectivamente en 1.0 y 2.1.

\medskip
\noindent
2.2.3. Supongamos ahora que $S$ es henseliano. Todo esquema cuasifinito
separado $X$ sobre $S$ posee entonces una descomposición canónica como
unión disjunta
\begin{equation}\tag{2.2.3.1}
 X=X^f\coprod X',
\end{equation}
donde $X^f$ es \underline{finito sobre $S$} y $X'_o=\varnothing$, es
decir, $X'$ está confinado a su fibra genérica $X'_\eta$; así, $X^f$
posee la misma fibra especial que $X$. Esta descomposición es
manifiestamente funtorial en $X$. En particular, si $X$ es un esquema
en grupos sobre $S$, entonces $X^f$ es un subesquema en grupos. Si $X$
es plano, respectivamente étale, sobre $S$, lo mismo ocurre con $X^f$.

Aplicamos esto a las componentes ${}_{\ell^\nu}A^o$ de $T_\ell(A^o)$.
Se obtiene un subsistema proyectivo de grupos
\begin{equation}\tag{2.2.3.2}
 \bigl(({}_{\ell^\nu}A^o)^f\bigr)_{\nu>0}
 \underset{\mathrm{dfn}}{=}T_\ell(A^o)^f
 \subset T_\ell(A^o),
\end{equation}
llamado la \underline{parte fija} del módulo de Tate $T_\ell(A^o)$.
Por construcción existen isomorfismos canónicos, funtoriales en $A$,
\begin{equation}\tag{2.2.3.3}
 T_\ell(A^o)^f\times_S s\xrightarrow{\ \sim\ }
 T_\ell(A^o)\times_Ss\xrightarrow{\ \sim\ }T_\ell(A_o^o).
\end{equation}
Por otra parte, la fibra genérica de $T_\ell(A^o)^f$ define un
subprogrupo canónico de $T_\ell(A_K)$, llamado también la
«\underline{parte fija}» de este módulo de Tate:
\begin{equation}\tag{2.2.3.4}
 T_\ell(A_K)^f\subset T_\ell(A_K).
\end{equation}

\medskip
\noindent
2.2.4. Cuando $\ell\ne p=\operatorname{char}k$, la parte fija
$T_\ell(A^o)^f$ puede interpretarse como un haz $\ell$-ádico constante
torcido sobre $S$ (pues los $({}_{\ell^\nu}A^o)^f$ son finitos étale
sobre $S$). Puesto que el funtor
\[
 X\longmapsto X_o=X\times_Ss
\]
de los esquemas finitos étale sobre $S$ a los esquemas finitos étale
sobre $k$ es una equivalencia de categorías cuando $S$ es henseliano,
existe una equivalencia análoga para las categorías de haces
$\ell$-ádicos constantes torcidos sobre $S$ y sobre $k$. Por tanto,
(2.2.3.3), que da la fibra especial de $T_\ell(A^o)^f$, caracteriza
esta parte fija salvo isomorfismo único. Por otra parte, como subhaz
$\ell$-ádico de $T_\ell(A^o)$, queda manifiestamente caracterizada por
su fibra genérica (2.2.3.4), o equivalentemente por el
$\mathbb Z_\ell$-submódulo correspondiente, estable bajo
$\pi=\operatorname{Gal}(\overline K/K)$:
\begin{equation}\tag{2.2.4.1}
 T_\ell(A_K)^f(\overline K)
 \underset{\mathrm{dfn}}{=}T_\ell(A(\overline K))^f
 \subset T_\ell(A(\overline K)).
\end{equation}
De hecho, se tiene el resultado siguiente.

\medskip
\noindent
\underline{Proposición} 2.2.5. \underline{Supongamos, como antes, que
$\ell\ne p$ y que $S$ es henseliano. Entonces el submódulo (2.2.4.1)
de $T_\ell(A(\overline K))=T_\ell(A_K)(\overline K)$ correspondiente a
la parte fija $T_\ell(A_K)^f$ de (2.2.3.4) es el submódulo
$T_\ell(A(\overline K))^I$ de invariantes bajo el grupo de inercia
$I\subset\pi$.}

Esto se deduce, pasando al límite inverso, de la relación análoga
válida para todo entero $m>0$ primo con
$p=\operatorname{char}k$:
\begin{equation}\tag{2.2.5.1}
 ({}_mA)^f(\overline K)=({}_mA(\overline K))^I.
\end{equation}
Para demostrar esta relación, introduzcamos el localizado estricto
$\widetilde S$ de $S$, cuyo cuerpo de fracciones $\widetilde K$ es la
subextensión no ramificada máxima de $\overline K/K$. Puesto que
$({}_mA)^f$ es finito étale sobre $S$, sus puntos con valores en
$\overline K$ ya tienen valores en $\widetilde K$. Así, el miembro
izquierdo de (2.2.5.1) es
$\operatorname{Hom}_S(\widetilde S,({}_mA)^f)$, mientras que el miembro
derecho es, por definición de $\widetilde K$, igual a
${}_mA(\widetilde K)$. Por la definición de $({}_mA)^f$, y dado que
$\widetilde S$ es un límite inverso de $S$-esquemas finitos étale
$S_i$, el primero es también igual a
\[
 \operatorname{Hom}_S(\widetilde S,{}_mA)
 ={}_m\operatorname{Hom}_S(\widetilde S,A).
\]
Por la propiedad definitoria del modelo de Néron,
$\operatorname{Hom}_S(\widetilde S,A)=A(\widetilde K)$, lo que
demuestra \sisi{(6.2.5.1)}{(2.2.5.1)}.

Para deducir 2.2.5, observemos que (2.2.5.1) da
\begin{equation}\tag{*}
 T_\ell(A(\overline K))^I
 =\varprojlim_\nu({}_{\ell^\nu}A)^f(\overline K)
 \supset\varprojlim_\nu({}_{\ell^\nu}A^o)^f(\overline K).
\end{equation}
Si $(x_\nu)$ pertenece al término situado más a la derecha, entonces,
para todo $\nu\geq0$, $x_\nu$ es una sección de $A$ infinitamente
divisible por $\ell$. Su valor en $s$ es, por tanto, un elemento de
$A_o(k)$ infinitamente divisible por $\ell$ y, en consecuencia,
pertenece a $A_o^o(k)$. Esto demuestra la igualdad de los dos últimos
términos de (*) y concluye la demostración de 2.2.5.

\medskip
\noindent
\underline{Proposición} 2.2.6. \underline{Sea $C_K$ un segundo esquema
abeliano sobre $K$, con modelo de Néron $C$, y considérese un
homomorfismo $u_K:C_K\longrightarrow A_K$, que se prolonga por tanto a
un homomorfismo $u:C\longrightarrow A$. Supongamos que $u_K$ es una
isogenia. Entonces, para todo primo $\ell$,
$T_\ell(u_K):T_\ell(C_K)\longrightarrow T_\ell(A_K)$ induce una
isogenia de las partes fijas}
\begin{equation}\tag{2.2.6.1}
 T_\ell(C_K)^f\longrightarrow T_\ell(A_K)^f,
\end{equation}
\underline{y, análogamente, la aplicación inducida por
$u_o:C_o\longrightarrow A_o$}
\begin{equation}\tag{2.2.6.2}
 T_\ell(u_o^o):T_\ell(C_o^o)\longrightarrow T_\ell(A_o^o)
\end{equation}
\underline{es una isogenia. Por último, $u_o:C_o\longrightarrow A_o$
induce una isogenia del mayor subtoro de $C_o$ al mayor subtoro $T_o$
de $A_o$; y sobre una extensión $k'/k$ en la que estén definidas las
«partes abelianas» de $C_o$ y $A_o$ (por ejemplo, sobre la clausura
perfecta de $k$), $u_o$ induce una isogenia entre las partes abelianas
de $C_{ok'}$ y $A_{ok'}$.}

Este enunciado se deduce formalmente de que $T_\ell(A_K)^f$,
$T_\ell(A_o^o)$, el mayor subtoro de $A_o$ y la «parte abeliana» de
$A_{ok'}$ son funtores aditivos del $A_K$ variable, y de que $u_K$ es
una isogenia precisamente en el sentido de que existe un morfismo
$v_K:A_K\longrightarrow C_K$ tal que
\[
 u_Kv_K=n\,\mathrm{id},\qquad v_Ku_K=n\,\mathrm{id}
\]
para algún entero $n>0$.

\medskip
\noindent
\underline{Corolario} 2.2.7. \underline{Si $A_K$ y $C_K$ son esquemas
abelianos isógenos sobre $K$, entonces los invariantes correspondientes
$\mu,\lambda,\alpha$ de (2.1.11) son los mismos para $A_K$ y $C_K$. En
particular, son los mismos para $A_K$ y su dual $A'_K$.}

\medskip
\noindent
\underline{Observaciones} 2.2.8. En los dos enunciados precedentes,
2.2.6 y 2.2.7, hemos supuesto que $S$ es henseliano. Sin embargo, es
inmediato que 2.2.7, así como las afirmaciones de 2.2.6 relativas a las
propiedades del homomorfismo inducido
$u_o:C_o\longrightarrow A_o$, son válidas sin esta hipótesis. En
efecto, se reduce al caso ya tratado mediante el cambio de base
$S^h\longrightarrow S$, donde $S^h$ es la henselización de $S$,
utilizando el hecho ---inmediato a partir de las definiciones--- de que
la formación de los modelos de Néron conmuta con este cambio de base.

\medskip
\noindent
\underline{Corolario} 2.2.9. (\underline{Buena reducción.})
\underline{Sea $S$ un rasgo, sea $A_K$ un esquema abeliano sobre su
cuerpo de fracciones $K$, sea $A$ el modelo de Néron de $A_K$, y sea
$\ell$ un primo distinto de $\operatorname{char}k$. Son equivalentes
las condiciones siguientes.}

\smallskip
\noindent
(i) \underline{Existe un esquema abeliano sobre $S$ cuya fibra genérica
es isomorfa a $A_K$.}

\noindent
(ii) $A$ \underline{es un esquema abeliano}.

\noindent
(ii$^o$) $A^o$ \underline{es un esquema abeliano}.

\noindent
(iii) $A_o$ \underline{es un esquema abeliano}.

\noindent
(iii$^o$) $A_o^o$ \underline{es un esquema abeliano}, \underline{es
decir, con la notación de (2.1.11), $\lambda=\mu=0$}.

\noindent
(iv) $A$ \underline{es propio sobre $S$}.

\noindent
(iv$^o$) $A^o$ \underline{es propio sobre $S$}.

\noindent
(v) $T_\ell(A_K)$ \underline{es «no ramificado sobre $S$»}, es decir,
\underline{se prolonga a un haz $\ell$-ádico constante torcido sobre
$S$}, o equivalentemente \underline{el grupo de inercia $I$ actúa
trivialmente sobre $T_\ell(A_K(\overline K))$}.

\noindent
(v bis) \underline{Para todo entero $\nu\geq0$,
${}_{\ell^\nu}A_K$ es «no ramificado sobre $S$»}, es decir,
\underline{se prolonga a un recubrimiento étale de $S$}, o
equivalentemente \underline{el grupo de inercia $I$ actúa trivialmente
sobre ${}_{\ell^\nu}A_K(\overline K)$}.

\smallskip
\noindent
\underline{Por último, un esquema abeliano sobre $S$ como en (i) es un
modelo de Néron de $A_K$ y, en particular, queda determinado salvo
isomorfismo único.}

\medskip
\noindent
2.2.9.1. Se dice que $A_K$ posee \underline{buena reducción} sobre $S$
si satisface la condición (i) anterior. El criterio (v), o (v bis), se
conoce como el \underline{criterio de Néron--Ogg--Shafarevich} de buena
reducción.

Demostremos 2.2.9. Es bien sabido (cf., por ejemplo, [5]) que un esquema
abeliano sobre $S$ satisface la propiedad universal de un modelo de
Néron; esto se reduce al hecho de que una sección racional de un esquema
abeliano sobre una base regular está definida en todas partes. Esto
demuestra la equivalencia
$(i)\Longleftrightarrow(ii)\Longleftrightarrow(ii^o)$, así como la última
afirmación de 2.2.9. Estas condiciones implican manifiestamente (iii),
(iv) y (iv$^o$), cada una de las cuales implica (iii$^o$). Demostremos
que (iii$^o$)$\Longrightarrow$(ii$^o$), lo que establecerá la
equivalencia de (i) a (iv$^o$). Esto se deduce del hecho de que un
esquema en grupos liso sobre $S$ cuyas fibras son esquemas abelianos es
propio sobre $S$ (y, por tanto, es un esquema abeliano), por
EGA~IV~15.7.10.

Por otra parte, en vista de 2.2.5 y del hecho de que
\[
 T_\ell(A_K(\overline K))/T_\ell(A_K(\overline K))^I
\]
es manifiestamente un $\mathbb Z_\ell$-módulo sin torsión, la condición
(v) expresa que $T_\ell(A_K(\overline K))$ y su parte fija tienen el
mismo rango. Por (2.1.11) y (2.2.3.3), esta es la condición
\[
 \mu+2\alpha=2(\lambda+\mu+\alpha),
\]
es decir, $2\lambda+\mu=0$, o $\lambda=\mu=0$, que es precisamente
(iii$^o$). Finalmente, la equivalencia entre (v) y (v bis) resulta
clara por las definiciones. Esto concluye la demostración de 2.2.9.
